\section{The \coolname{} Dataset}\label{sec:dataset}

The methodology of \S\ref{sec:method} requires a dataset with four properties:
(i)~the same physical scenes captured across controlled phase transitions,
(ii)~multiple perception modalities with shared ground truth,
(iii)~6-DoF camera pose for temporal metrics, and
(iv)~objects described in natural language for grounding evaluation.
We are not aware of any existing dataset that provides all four (Table~\ref{tab:dataset-comparison}).
We construct \coolname{} to fill this gap.

\paragraph{Sensor platform.}
A hand-carried rig pairs an Intel RealSense \textbf{D435} (RGB-D, structured IR, 0.3--3\,m) with a \textbf{T265} (stereo visual-inertial, 6-DoF pose at 200\,Hz) and a head-mounted \textbf{GoPro} during interaction.
Human-carried capture is a deliberate proxy for embodied robot motion, validated via motion-profile comparison with deployed platforms (Appendix~\ref{sec:appendix-motion}):
median $\Delta t = \todo{X}$\,cm/frame, median $\Delta R = \todo{X}$$^\circ$/frame, within the envelope of Stretch RE2, Unitree H1, and Fetch.

\paragraph{Scale and diversity.}
99 real-world scenes (\todo{X} indoor, \todo{X} outdoor), spanning daylight and night, reflective and transparent surfaces.\footnote{The collection target was 100 scenes; one scene was excluded from v1 due to an incomplete capture and is planned for v1.1.}
${\sim}$70 daily object categories (\todo{X} total physical instances) following FewSOL~\citep{padalunkal2023fewsol}: real household items varying in material, surface, and size.
${\sim}$75{,}000 evaluation frames across three phases at $640 \times 480$, 30\,fps.
\textbf{Modalities:} RGB, metric depth, 6-DoF camera pose, instance masks, tracklets, per-object textual attributes, T265 fisheye stereo pairs, and GoPro egocentric video during interaction.

\paragraph{Objects and language annotations.}
Per-object textual attributes (category, colour, material, shape) use layperson vocabulary collected via Amazon Mechanical Turk for the FewSOL object set---matching how real users will command robots.
The same objects appear in both single-object (isolated, AR-tagged) and multi-object (cluttered, three-phase) scenes, enabling in-context visual prompting evaluation (\S\ref{sec:tasks}).

\paragraph{Ground truth.}
\textit{Instance masks:} SAM2 auto-mode + GroundingDINO proposals + iterative human refinement; refinement count $k_j$ feeds ESD (\S\ref{sec:features}).
\textit{Tracklets:} SAM2 video propagation with human-corrected keyframes.
\textit{Depth:} raw D435 structured-light; invalid pixels flagged, no inpainting.
\textit{Camera pose:} T265 VIO at 200\,Hz, downsampled to 30\,fps.
Pose quality is reliable during clutter and clean phases (smooth walkaround motion) but degrades during interaction (hand occlusion of T265 IR emitters).
TS is therefore restricted to non-interaction phases in this release (\S\ref{sec:deploy-metrics}).
\textit{QA ground truth:} human-authored question templates instantiated deterministically from annotations (\S\ref{sec:tasks}).
Full annotation pipeline details and figures in Appendix~\ref{sec:appendix-annotation}.

\paragraph{Availability.}
HuggingFace (\url{\webpage}), CC BY 4.0.
5-scene reviewer sample ($<$4\,GB); full dataset ${\sim}$\todo{X}\,GB.
Croissant metadata (core + RAI fields) included.
All 99 scenes are released for evaluation; \coolname{} is a benchmark for models trained elsewhere, so no train/val split is provided.
Difficulty stratification (Easy/Medium/Hard) is shipped at two granularities in the splits manifest:
(i)~\texttt{scene\_splits.json}---a per-scene 33/33/33 (33/33/34 once the 100\textsuperscript{th} scene lands) tier assignment, the primary deliverable for benchmark consumers, with all three phases of a scene held in the same tier to satisfy the state-transition robustness requirement;
(ii)~\texttt{phase\_difficulty.json}---per-(scene, phase) structured metric with the primary effort\_stratified\_v1 tier, cross-method consensus tier, per-row weight-perturbation confidence flag, and an 8-d per-modality sub-score vector that preserves the multi-faceted nature of difficulty (\S\ref{sec:features}, Appendix~\ref{sec:appendix-method-study}).
